% Bagian: model-theory
% Bab: interpolation
% Subbagian: introduction

\documentclass[../../../include/open-logic-section]{subfiles}

\begin{document}

\olfileid[id]{mod}{int}{int}

\olsection{Pendahuluan}

Teorema interpolasi adalah hasil berikut: Andaikan
$\Entails !A \lif !B$. Maka terdapat !!a{sentence} $!C$ sedemikian sehingga
$\Entails !A \lif !C$ dan $\Entails !C \lif !B$. Selain itu, setiap
!!{constant}, !!{function}, dan !!{predicate} (selain $\eq$) dalam
$!C$ terdapat dalam $!A$ maupun~$!B$. !!^{sentence} $!C$ disebut
\emph{interpolan} dari $!A$ dan~$!B$.

Teorema interpolasi menarik dalam dirinya sendiri, tetapi kepentingan
utamanya terletak pada kenyataan bahwa teorema ini dapat digunakan untuk
membuktikan hasil-hasil tentang keterdefinisian dalam suatu teori dan tentang
syarat-syarat yang membuat penggabungan dua teori konsisten menghasilkan teori
yang konsisten. Hasil pertama dikenal sebagai teorema keterdefinisian Beth;
hasil kedua sebagai teorema konsistensi gabungan Robinson.

\end{document}
